\documentclass{article}
\usepackage[UTF8]{ctex}
\usepackage{amsmath,amssymb,amsfonts}
\usepackage{enumitem}
\usepackage{geometry}
\geometry{a4paper, margin=1in}
% === 错题本排版优化（不参与题目标识） ===
\setlength{\parindent}{0pt}
\setlength{\parskip}{0.35em plus 0.10em minus 0.08em}
\setlength{\emergencystretch}{3em}
\setlist{itemsep=0.25em, topsep=0.40em, parsep=0pt}
\allowdisplaybreaks[2]
% === 错题本排版优化结束 ===
\begin{document}
\section*{中国科学技术大学2026年数学分析真题}
\section*{106．中国科学技术大学 2026 年数学分析试卷}

\subsection{题目1}
1．求极限 \(\lim _{n \rightarrow \infty} \sum_{k=1}^{n} \frac{1}{n^{2}+k}\) ．

\subsection{题目2}
2．设 \(x_{0} \in\left(0, \frac{\pi}{2}\right)\) ，且 \(x_{n}=\sin x_{n-1}(n=1,2, \cdots)\) ，求 \(\lim _{n \rightarrow \infty} x_{n}, \lim _{n \rightarrow \infty} n x_{n}^{2}\) ．

\subsection{题目3}
3．设 \(f(x)=\ln \frac{1+x}{1-x}\) ，求 \(f^{(n)}(0)\) ．

\subsection{题目4}
4．设 \(y=f(x)\) 是 \(e^{x}+x y=e\) 在 \((1,0)\) 点附近确定的隐函数，求 \(f^{\prime \prime}(0)\) ．

\subsection{题目5}
5．求 \(z=\frac{1}{2}\left(x^{2}+y^{2}\right)\) 被柱面 \(\left(x^{2}+y^{2}\right)^{2}=2 x y\) 所截下的面积．

二．（15 分）设函数 \(f(x, y)= \begin{cases}\frac{|x y|^{\alpha}}{x^{2}+y^{2}}, & (x, y) \neq(0,0) ; \\ 0, & (x, y)=(0,0) .\end{cases}\)

\begin{itemize}
  \item[（1）] \(\alpha\) 取何值时，\(f(x, y)\) 在 \((0,0)\) 处连续？
  \item[（2）] \(\alpha\) 取何值时，\(f(x, y)\) 在 \((0,0)\) 处可微？
\end{itemize}

三、15 分）设 \(f(x)\) 在 \(\mathbb{R}\) 上反常可积，即 \(\int_{-\infty}^{+\infty}|f(x)| \mathrm{d} x\) 收敛，定义函数 \(\varphi(x)=\int_{-\infty}^{+\infty} f(t) \cos (x t) \mathrm{d} t\) ，证明：\(\varphi(x)\) 在 \(\mathbb{R}\) 上一致连续。

四．（20 分）设 \(\Omega\) 为 \(\mathbb{R}^{3}\) 中一有光滑边界的有界区域，设 \(u \in C^{2}(\Omega) \cap C^{1}(\bar{\Omega})\) ，且 \(\Delta u=0\) 。

\begin{itemize}
  \item[（1）] 证明： \(\iint_{\partial \Omega} u \frac{\partial u}{\partial n} \mathrm{~d} S=\iiint_{\Omega}\|\nabla u\|^{2} \mathrm{~d} V\) ．
  \item[（2）] 若 \(\left.u\right|_{\partial \Omega} \equiv 1\) ，证明：\(u \equiv 1\) ．
\end{itemize}

五．（15 分）证明：

\begin{itemize}
  \item[（1）] 对任意的 \(x \in \mathbb{R}\) ，有 \(|\cos x|=\frac{2}{\pi}+\frac{4}{\pi} \sum_{n=1}^{\infty} \frac{(-1)^{n+1}}{4 n^{2}-1} \cos (2 n x)\) ．
  \item[（2）] \(\lim _{\lambda \rightarrow+\infty} \int_{0}^{1} f(x)|\cos \lambda x| \mathrm{d} x=\frac{2}{\pi} \int_{0}^{1} f(x) \mathrm{d} x\) ，其中 \(f(x)\) 在［0，1］上黎曼可积。
\end{itemize}

六．（15 分）设 \(f(x)\) 在 \([0,+\infty)\) 上可导，且 \(0 \leq f(x) \leq \frac{x}{1+x^{2}}\) ，证明：存在 \(\xi \in(0,+\infty)\) ，满足

\[
f^{\prime}(\xi)=\frac{1-\xi^{2}}{\left(1+\xi^{2}\right)^{2}}
\]

七．（10 分）设 \(f(x) \in C[0,1]\) ，对任意的 \(0 \leq a<b \leq 1, f\) 在 \([a, b]\) 上至少有两个不同的最大值点，证明：\(f(x)\) 为常值函数．

八．（10 分）设 \(f:[0,1] \rightarrow[0,1]\) 满足

\[
|f(x)-f(y)|<|x-y|, \forall x, y \in[0,1], x \neq y .
\]

证明：存在唯一的 \(\xi \in[0,1]\) ，使得 \(f(\xi)=\xi\) ．
\end{document}
